\section{Limitations}
\label{sec:limitations}

The principal limitation of this study is construct validity. The model
observes adjudicated scores, caption-level score components, timing, class
membership, competition context, and prior score history. It does not observe
the performed show itself, nor most factors that shape it: musical or visual
execution, artistry, instruction quality, performer membership, show design,
rehearsal time, resource availability, or how fully the show had been
staged and learned at debut. As a result, the model estimates continuity within the
adjudication system rather than any feature of the underlying activity. A low
prediction error means that debut adjudication agrees with championship
adjudication; it does not mean the model explains why an ensemble improved,
whether scores accurately reflect artistic quality, or how any ensemble should
train or be coached. MAE measures agreement with later judges, not accuracy
against an objective ground truth of performance quality.

Additional limitations include: championship-participant selection bias (only
ensembles that reached a championship contribute a target); unequal schedules
and sparse independent-class cells (PIO and PIA have two and six held-out
observations, respectively); class changes and incomplete knowledge of
intended class at debut (adjudicators may promote ensembles mid-season);
debut show completeness (how fully the production has been staged at the time
of the first performance), which is not recorded and may vary more across
A-class ensembles than in World class, a plausible partial explanation for
PSA's wider error distributions that cannot be tested with the available data;
pandemic discontinuity and uncertain historical sheet revisions; and
SCPA-only geographic scope, which limits generalizability to other circuits.

Judge assignments in this dataset are observational and cannot be assumed to
be randomized. Because the same judges work multiple events across a regional season, scores from repeated judge--ensemble pairings are not
independent. Controlling for confounders would require assignment data not
available here; exploratory residual screening identified patterns that could
not be attributed defensibly to individual judges under these constraints, and
those findings are not presented as results.

Prediction intervals are marginal empirical bands derived from development
out-of-fold residuals, not calibrated ensemble-specific guarantees.
